% !PRINTED 472
\weberpage{472}

But equation (22) again has as its consequence an equation of the form (21), so that the signs agree in both, and equation (21) requires, by § 40, (7):
\weq{23}{
\begin{aligned}
&\frac{c + \partial\omega}{a} = \frac{\gamma + \delta\omega}{\alpha + \beta\omega}, \quad \delta \equiv 0\ (\mathrm{mod}\ 2), \\
&\gamma^2 - 1 + \gamma(2\alpha - \delta) \equiv 0 \ \text{or} \equiv 8 \quad (\mathrm{mod}\ 16),
\end{aligned}
}
and in fact, according as the one or the other congruence holds, the one or the other sign applies in (21) and (22).

If we take $m = n$, and therefore $m \equiv -1\ (\mathrm{mod}\ 8)$, then from (2) and (23) there results, exactly as above:
\[
\begin{aligned}
\beta &= A, & n\gamma &= -C + cB, \\
\alpha &= B, & n\delta &= -B + cA,
\end{aligned}
\]
and hence, from the second congruence (23):
\[
C^2 + CB - 1 \equiv 0 \ \text{or} \equiv 8 \quad (\mathrm{mod}\ 16),
\]
from which it follows first of all that $B$ must in any case be divisible by 8; but if we increase $B$ by an odd multiple of 8, then one of these congruences changes into the other, and consequently one sign in (21) changes into the other. For a definite $\omega$, therefore, only one of the two equations (22) holds, and it follows:

7. \gesperrt{If $m \equiv 7\ (\mathrm{mod}\ 8)$, then $\sqrt{2}\, f(\omega)^3$ is a class invariant.}

If we may trust an induction confirmed by all known examples, then the following theorem also holds:

8. \gesperrt{If $m \equiv 3\ (\mathrm{mod}\ 8)$, then $f(\omega)^3$ is a class invariant.}

The general proof of this is, however, still lacking.

If $m \equiv 2\ (\mathrm{mod}\ 4)$, we apply the same procedure as above in the case $m \equiv 1\ (\mathrm{mod}\ 8)$.

We put:
\weq{24}{2n = m + y^2,}
and take $y = 0$ or $= \pm 2$, so that $n \equiv 1\ (\mathrm{mod}\ 4)$ results. % [sic] German print omits y before the second equality

If we then set in equation (19)
\weq{25}{v_1 = f_2\left(\frac{\omega}{2}\right)^3 = \frac{\sqrt{2}^{3}}{f_1(\omega)^3}, \quad u_1 = f_1(\omega)^3}
then there results an equation which contains only $\sqrt{2}\, f_1(\omega)^6$ and rational coefficients. (19), however, is satisfied only [§ 40, (7)] if
%CONT>>
